\section{Deterministic implementation and completion}
\label{sec:derandomization}

The preceding constructions yield an additive approximation in expectation.
We now replace their random choices by finite distributions of polynomial
support, and then choose the least expensive feasible outcome. Throughout
this section the desired approximation loss is fixed. In particular, the
grid denominator, the number of physical rounds, and all parameters chosen
in Theorems~\ref{thm:preparation} and~\ref{thm:grid-rounding} are constants
independent of the input size. All input distances are rational, and
running time includes their encoding lengths.

\subsection{Compression of finite distributions}

A distribution list is a finite list of states, each equipped with a
nonnegative rational weight, whose weights sum to one. The state records
the data needed to continue a construction; a feature vector records the
statistics whose expectations must be retained.

The next lemma is the constructive finite-dimensional Carath\'eodory
argument for preserving an expectation. Its connection with moment
matching is classical~\cite{BayerTeichmann2005}; we include the exact
elimination and its bit bound because both are needed by the algorithm.

\begin{lemma}[Affine support reduction]
\label{lem:affine-support}
Let a distribution list have states \(s_1,\ldots,s_M\), weights
\(\lambda_1,\ldots,\lambda_M\), and feature vectors
\(v(s_i)\in\mathbb Q^d\). There is an exact algorithm that replaces the
weights by nonnegative rational weights supported on at most \(d+1\) of
the same states, preserving
\[
  \sum_{i=1}^M\lambda_i v(s_i).
\]
Its running time and output encoding length are polynomial in the
encoding length of the given list and feature vectors.
\end{lemma}

\begin{proof}
Discard zero-weight entries. If more than \(d+1\) remain, choose a
nonzero integral affine dependence \(z\) satisfying
\[
 \sum_i z_i=0,
 \qquad
 \sum_i z_i v(s_i)=0.
\]
Since its coefficients sum to zero, there are both positive and negative
coefficients. Set
\[
 t=\min_{i:z_i>0}\frac{\lambda_i}{z_i},
 \qquad
 \lambda'_i=\lambda_i-tz_i.
\]
The new weights are nonnegative, have the same sum and feature
expectations, and include at least one additional zero. Repeating gives
the support bound.

For completeness, the rational arithmetic has polynomial complexity.
An affine dependence can be found using at most \(d+2\) columns of the
matrix with columns \((1,v(s_i))\). Clearing denominators and using
minors gives an integral dependence of polynomial bit length. If
\(\lambda_i=b_i/D\) for a common denominator \(D\), and \(j\) attains
the minimum defining \(t\), then
\begin{equation}
 \lambda'_i=
 \frac{b_i z_j-b_jz_i}{Dz_j}.
 \label{eq:compression-denominator}
\end{equation}
Thus one elimination adds at most the bit length of \(|z_j|\) to a
common denominator. There are at most \(M\) eliminations, and the
feature vectors do not change. The resulting bit lengths and the exact
linear-algebra computations are polynomially bounded.
\end{proof}

Support reduction preserves every property satisfied by each state in
the list, whether or not that property appears among the features. It
does not preserve the entire joint distribution or arbitrary future
adaptive behavior. We therefore specify the retained features at each
use and verify the required estimates for the resulting process.

\subsection{A distribution list for the preparation stage}

Let \(M_\star\) denote the number of stars in the preparation
construction, and write \(X_e\) for the current count of star \(e\).
Each count remains within its original floor and ceiling. A weighted
pair update within a strict sign preserves the weighted sum and has
conditional mean equal to the current count vector. For every pair of
distinct coordinates \(e,f\), it also satisfies
\begin{equation}
 \E[X'_e\mid X]=X_e,
 \qquad
 \E[X'_eX'_f\mid X]\le X_eX_f.
 \label{eq:star-local-moments}
\end{equation}
Indeed, unless both coordinates are the selected pair, the product is
affine in the update parameter. If they are the selected pair, their
increments have opposite signs and the product is concave. A
one-coordinate rounding with the correct mean preserves every distinct
pair-product expectation.

These assertions hold for every admissible state and every permissible
choice of the next pair. Thus the next pair may be selected adaptively,
for example as the first two fractional coordinates in a fixed order.
After branching one update at every state, apply
Lemma~\ref{lem:affine-support} with the features
\[
   X_e \quad\text{and}\quad X_eX_f \qquad(e<f).
\]
Induction using~\eqref{eq:star-local-moments} preserves the initial
first moments and the upper bounds on distinct pair moments. Full
independence between the positive, negative, and zero-coefficient
classes is unnecessary for these conclusions.

The weighted sums and floor/ceiling restrictions are properties of each
retained state. Each productive pair update makes at least one
fractional count integral, and integral counts are never updated again.
There are therefore only \(O(M_\star)\) updates, with identity updates
padding states that finish earlier. Ordinary one-coordinate rounding
also handles the zero coefficients and the final positive-side
fraction. We leave the final negative-side fraction unrounded until
these operations are complete. Each compressed list has
\(O(M_\star^2)\) entries.

\paragraph{The final negative-side fraction.}
Condition on a state of this last list. There is at most one fractional
count, say \(X_e=m+\xi\), where \(0<\xi<1\). The identity of \(e\)
may depend on the state. Perform the special rounding of the
preparation construction: use a nominal coin of mean \(\xi\), open the
outside supplier according to \(m+1\) if its flag is one, and activate
the extra deficit on a subset of \(D_e\) with marginal \(\xi\) at
every target.

There is a polynomial-support law for this subset. If
\(d_e=|D_e|\), choose a size
\[
 L\in\{\lfloor\xi d_e\rfloor,\lceil\xi d_e\rceil\},
 \qquad \E L=\xi d_e,
\]
then choose a uniformly random cyclic interval of length \(L\) in a
fixed cyclic order of the targets. Each target has inclusion
probability \(\E L/d_e=\xi\), and the support has size at most
\(2d_e\). Including the nominal coin changes this bound by only a
constant factor.

For a fixed target, the special star's activation count has conditional
mean \(X_e\); its nominal count has the same conditional mean. Every
other star count is integral in the conditioned state. Consequently,
for any two distinct stars activating this target, the conditional
expectation of their activation product is the product of their
current counts. This argument applies separately in every state, so
it remains valid when the identity of the final fractional star varies.
After averaging, the activation counts retain the nonpositive
pair-covariance bounds. Each individual activation count remains
between the original adjacent integers, giving its variance bound as
well. No pairwise independence between different targets of this last
star is required: the deficit estimates are made for one target at a
time, and the clipping bound sums expectations over targets.

\paragraph{Conditional rounding inside bundles.}
For each resulting outer state, clip the deficits and rescale the
bundle masses exactly as in Theorem~\ref{thm:preparation}. Keep this
outer state fixed while constructing a new list for all the
within-bundle rounds. Successive pair updates and compression of all
first and distinct-pair moments of the opening vector preserve its
marginals and nonpositive covariances conditional on this outer
state. Updates in different bundles preserve cross-bundle product
expectations; thus their entire Cartesian product need not be
enumerated.

In particular, if \(\eta_i\) is the rounded opening minus its
conditional mean in the fixed outer state, then the noise on a
client's old assignment obeys
\[
 \Var\left(
   \sum_{i\in F_j\cap\bigcup_sP_s}\eta_i
   \,\middle|\,\text{outer state}
 \right)
 \le
 g\sum_s\sum_{i\in F_j\cap P_s}\frac{p'_s}{p_s}y_i
 \le g/p_0\le 2g.
\]
Here \(p'_s\le1\) and \(p_s\ge p_0\ge1/2\). This is the conditional
noise estimate used in the preparation proof. Maintaining a separate
inner list for every outer state retains this estimate despite
the varying rescalings.

All remaining preparation estimates use the retained first and second
moments, or properties holding in every state. In particular, supplier
availability and the raw opening-count bound hold in each state;
the clipping estimate and the probability of exceeding its count
allowance follow from the same deficit moments. The preparation list
therefore satisfies the same unconditioned cost and budget-success
bounds as the randomized construction. Discarding budget failures
and selecting a remaining grid solution of minimum total nearest-unit
assignment cost gives the corresponding deterministic bound on the
sum over clients. This single choice need not simultaneously attain
each client's individual expectation bound.

\subsection{A distribution list for each physical round}

Fix a physical state and a round. For each target copy, first construct
a small list of source sets of cardinality \(\Delta\) with the
prescribed edge marginals. Use ordinary fixed-sum pair rounding:
after each pair update, branch each current state to its at most two
endpoints and immediately compress the resulting list while preserving
first moments. Thus no full tree of a target's pair updates is ever
enumerated, and every intermediate list has polynomial support.
All sets in this list
contain the target's own group and respect the edges whose prescribed
marginal is zero or one.

Combine the target lists by successive independent products. After
each product, compress while retaining every edge first moment and
every product of two edge indicators belonging to distinct targets.
On adjoining a new target independently, its edge products with earlier
targets have the required expectations; compression then preserves
them. Induction gives a graph list with all the first and cross-target
second moments of the independent-target graph construction. Its
feature dimension and support size are polynomial in the number of
active copies.

The ideal generator in~\eqref{eq:rnd-ideal-generator} is linear in
edge indicators except for the terms \(\delta_hY_h\). Expanding these
terms gives products of edges from a common source to distinct targets,
and diagonal terms reducible to first moments. The retained statistics
therefore suffice for its graph average. The heavy-source rule and
forced openings depend on more than these moments, but this introduces
no extra averaging requirement: the batch drift
bound~\eqref{eq:rnd-batch-drift} holds conditional on every realized
graph, with that graph's own forced openings and chunks. Only after
that bound is applied do we average the ideal generator, obtaining
\eqref{eq:rnd-average-drift} for the compressed graph distribution.

\paragraph{Proposal and trigger lists.}
For every graph, include both round types, each with probability
\(1/2\). Within either type, construct the proposal or trigger
indicators by successive independent Bernoulli products, compressing
after each product while retaining all squarefree monomials of degree
at most four. These moments coincide exactly with those of the
independent reference law. There are polynomially many such
monomials. In the balanced type the variables being compressed are
the proposals, before conflict filtering.

To see why filtering requires no higher moments, write \(Q_h\) for a
balanced proposal and \(I_h\) for its acceptance. If \(C_h\) is the
set of sources conflicting with \(h\), then, in every outcome,
\[
 Q_h-\sum_{t\in C_h}Q_hQ_t\le I_h\le Q_h,
 \qquad
 I_hI_t\le Q_hQ_t\quad(h\ne t).
\]
A balanced source has \(|E_h|=\Delta\), and each target has
\(\Delta\) possible sources, so \(|C_h|\le\Delta(\Delta-1)\).
Preservation of first and second proposal moments gives
\begin{equation}
 p_*-\Delta(\Delta-1)p_*^2\le\E I_h\le p_*,
 \qquad
 \E[I_hI_t]\le p_*^2\quad(h\ne t).
 \label{eq:list-conflict-moments}
\end{equation}
These are the marginal and pair bounds needed for the isolated-action
and simultaneous-action estimates. Accepted out-sets remain disjoint
in every outcome.

For an unbalanced round, the count statistic from
Equation~\eqref{eq:rnd-count-statistic}, conditional on the graph,
has the form
\[
 Z_*=
 \sum_e c_e I_e+\sum_{e<f}c_{ef}I_eI_f
\]
with fixed rational coefficients. It has degree at most two in the
trigger indicators, and its square has degree at most four after using
\(I_e^2=I_e\). The compressed law therefore has exactly the same
\(\E Z_*\), \(\E Z_*^2\), and \(\Var Z_*\) as the independent
reference law. The variance estimate proved for independent
indicators transfers to this law, and hence so does the Chebyshev
count-failure bound. An independent-coordinate variance inequality
is applied to the reference law before this transfer.

The terminal alive-mass estimate likewise uses only the marginal and
pair bounds for the \(\Delta\) possible incoming triggers at each
surviving target. It follows by the same union lower bound as in the
randomized analysis. Thus all probabilistic estimates in
Theorem~\ref{thm:grid-rounding} remain valid for these lists.

\paragraph{Retaining adaptive histories.}
Expand these graph and proposal lists at every reached physical
history, and retain all resulting physical outcomes over the
\(T\) rounds. The history tree is not compressed. The alive sets
used in the analysis can be defined along this tree, but need not be
computed by the algorithm. Even when two histories reach the same
physical state with different analysis-only alive sets, the
next-round estimates apply to either history, since they hold for
every admissible alive set conditional on the current physical state
and graph.

The number of active copies is polynomial: there are at most
\(\Delta\) at each facility, and resetting regenerates copies at an
existing facility. The number of chunk actions is also polynomial,
since the total number of source--target incidences is
\(\Delta\) times the number of target copies. Each history has only
polynomially many children. Since \(T\) is an input-independent
constant, the complete history tree has polynomial size. At each
leaf perform the final deterministic completion, discard outcomes
exceeding the additive budget, and select one of minimum total
connection cost.

\subsection{Rational implementation and bit complexity}

We give the encoding details because support size alone does not imply
a polynomial-time algorithm. Fix rational values of all algorithmic
parameters. In each of the finitely many rounds, fix in advance a
rational taper value \(\widehat P\) satisfying its stipulated
accuracy. The analytic functions of elapsed time, the potential, and
the analysis-only alive-set decisions need not be evaluated by
the algorithm.

An optimal basic LP solution has polynomial encoding length. Splitting
facilities to make assignments use whole pieces introduces only
polynomially many pieces. Piece masses, bundle masses, and
intersection masses are sums and differences of the LP coordinates.
The extra star weights multiply a deficit by the fixed rational
\(\gamma\), and conversion to grid counts multiplies by the integer
\(\Delta\). Hence the initial star counts have a common denominator
\(D_0\) of polynomial bit length.

For weighted pair rounding, write \(a_e=|u_e|\). If a pair update
makes \(X_e\) integral and leaves \(X_f\) fractional, preservation
of its weighted sum gives
\[
 a_fX'_f=a_fX_f+a_eX_e-a_eX'_e.
\]
Throughout these updates, therefore, every weighted count lies in
\(D_0^{-1}\mathbb Z\). Recovering \(X_f\) only divides by the
polynomially bounded integer \(a_f\). Counts remain in their original
integer bounds, so their numerators and their endpoint probabilities
have polynomial bit length as well. Final single-coordinate rounding
makes the coordinate integral; zero-coefficient counts retain their
initial values until rounded directly. Clipping and rescaling give
fresh rational inputs of polynomial bit length for the within-bundle
rounds, where the same argument applies with unit coefficients.

The taper radii can also be computed exactly. On an interval where
\(m\) source copies have saturated marginal one and \(r\) source
copies have marginal strictly between zero and one, the radius
equation is
\[
 m+\frac{r}{\widehat P}
 -\frac{\sum_{h\text{ partial}}d(i,h)}{\widehat P N_i}
 =\Delta.
\]
When \(r>0\), an interior solution is
\begin{equation}
 N_i=
 \frac{\sum_{h\text{ partial}}d(i,h)}
      {r+\widehat P(m-\Delta)}.
 \label{eq:rational-taper-radius}
\end{equation}
The breakpoints are the rational numbers \(d(i,h)\) and
\(d(i,h)/(1-\widehat P)\). Sorting these values and checking the
successive intervals finds the smallest solution, including a
breakpoint solution. If \(r=0\), the equation is constant on the
interval, and the first point of a solution interval is a breakpoint.
The separately prescribed case of sufficient
co-located mass uses \(N_i=0\). All radii, marginals, and resulting
branch probabilities have polynomial bit length. The physical states
and their feature values are determined without reference to the
list weights.

It remains to control the list weights across repeated branching and
compression. Suppose a feature vector has \(d\) coordinates of bit
length at most \(L\). In the affine-dependence computation, clearing
denominators in each row of the matrix on at most \(d+2\) columns
produces integer entries of bit length \(O(dL)\). A nonzero kernel
vector can be chosen using minors, with bit length
\(O(d^2L+d\log(d+1))\) by the determinant bound. These bounds are
independent of the current list weights.

By~\eqref{eq:compression-denominator}, an elimination multiplies a
common weight denominator by just one coefficient of such a
dependence. An exact branch with probability denominator \(q\)
multiplies a common denominator by at most \(q\). Branching a whole
list adds the bit lengths of only polynomially many such
denominators, and taking a product of lists adds their denominator
bit lengths. There are polynomially many list operations, including
the entire constant-depth physical history tree. Their common
denominators and numerators consequently retain polynomial bit
lengths. Exact linear algebra, comparisons, and cost evaluations
therefore give a deterministic algorithm whose bit complexity is
polynomial in the input encoding length for every fixed choice of
accuracy parameters.

\subsection{Loss allocation and the approximation theorem}

We first record explicitly how conditioning and deterministic
selection are combined. For a nonnegative cost \(C\) and an event
\(G\) with \(\Prob(G)\ge1-r\),
\begin{equation}
 \E[C\mid G]
 \le\frac{\E C}{1-r}.
 \label{eq:conditioning-cost}
\end{equation}
Thus a finite distribution list contains a successful outcome whose
cost is no more than the right-hand side. No independence between
cost and budget success is required.

\begin{theorem}[Deterministic additive approximation]
\label{thm:pseudo}
For every fixed \(\varepsilon>0\), there is an
instance-independent integer \(c\ge0\) and a deterministic
polynomial-time algorithm for metric \(k\)-median that opens at
most \(k+c\) facilities and has total connection cost at most
\[
 \left(1+\frac2e+\varepsilon\right)\operatorname{LP},
\]
where \(\operatorname{LP}\) is the optimum value of the assignment
relaxation. In particular, its cost is at most the same factor times
\(\OPT_k\).
\end{theorem}

\begin{proof}
Put \(\alpha_\star=1+2/e\), and choose a fixed positive rational
number
\[
 t\le\min\{1/8,\varepsilon/10\}.
\]
Choose the parameters in the constructions underlying
Theorems~\ref{thm:preparation} and~\ref{thm:grid-rounding} with
unconditioned cost factors at most \(1+t\) and
\(\alpha_\star+t\), respectively, and with their respective
budget-failure probabilities at most \(t\). The proofs of those
theorems allow the cost and failure tolerances to be made
arbitrarily small; their stated expectation guarantees absorb this
same conditioning loss. Here we keep the loss visible.

All parameter choices can be made in a consistent order. After the
preparation parameters other than the grid are fixed, take an even
integer \(\Delta\) sufficiently large, set \(g=1/\Delta\), and
take a rational horizon of order
\(1+\log(\Delta/t)\). The requirements that the terminal error
be small and that \(\Lambda g\) be small are compatible because
\(\log\Delta/\Delta\to0\). Then choose the taper tolerance,
the selection bias, and the time step successively small. The
number of rounds is now fixed. Finally choose the heavy-source
parameter sufficiently small that the sum of the conditional
count-failure probabilities over all rounds is at most \(t\).
The resulting additive opening allowances are constants depending
only on these parameter choices.

Let \(L=\operatorname{LP}\). Apply the preparation list to an
optimal LP solution. Its total nearest-unit cost has expectation at
most \((1+t)L\), and its budget event has probability at least
\(1-t\). Select the least expensive successful grid solution.
By~\eqref{eq:conditioning-cost}, its total nearest-unit cost
\(L_{\rm grid}\) satisfies
\[
 L_{\rm grid}\le\frac{1+t}{1-t}L.
\]
Apply the physical-rounding list to this fixed grid solution. The
summed cost bound is at most \((\alpha_\star+t)L_{\rm grid}\)
before discarding count failures, and the budget succeeds with
probability at least \(1-t\). Selecting the least expensive
successful integral output therefore gives cost at most
\[
 \frac{(1+t)(\alpha_\star+t)}{(1-t)^2}L.
\]
Since \(\alpha_\star<2\) and \(t\le1/8\),
\begin{align*}
 \frac{(1+t)(\alpha_\star+t)}{(1-t)^2}-\alpha_\star
 &=
 \frac{(3\alpha_\star+1)t+(1-\alpha_\star)t^2}
      {(1-t)^2}\\
 &\le\frac{7t}{(1-t)^2}
 <10t
 \le\varepsilon.
\end{align*}
The two additive allowances sum to an instance-independent
constant, which may be rounded up to an integer \(c\).

Both minimum-cost selections are computable on polynomial-size
lists, with polynomial bit complexity as proved above. They use
total costs; the deterministic grid solution is not required to
attain every client's individual expectation bound. If \(L=0\),
the nonnegative cost has expectation zero in the preparation list
and then in the physical-rounding list, so the same argument gives
zero cost without division by \(L\).

Every opened copy lies at an original candidate facility. Merge
copies introduced by the construction back to their original
facility. This preserves connection distances and cannot increase
the number of opened facilities. Finally,
\(\operatorname{LP}\le\OPT_k\), proving the last assertion.
\end{proof}

\subsection{Applying the reduction on rational inputs}
\label{sec:reduction-implementation}

Before applying Theorem~\ref{thm:li-svensson}, we record its interface
with the deterministic algorithm just proved. Every call of that reduction
uses the same clients and distances with a restricted candidate set
\(F'\subseteq F\); our pseudo-algorithm uses no additional input promise
or prescribed-open facilities. Its internal forced openings are decisions
of that call, not constraints passed from another call. Copies are merged
back to original candidates before returning its output. If
\(0<|F'|<k\), opening every candidate gives the optimum for that restricted
instance. If \(F'=\varnothing\) and clients remain, discard the branch.
The distinguished branch retaining a padded optimal \(k\)-set has
neither exception. When the additive allowance is zero, the
pseudo-algorithm already satisfies the required budget.

The zero-optimum case can be recognized before the reduction. Every
client location must be an available candidate location, and there can
be at most \(k\) distinct client locations. Equivalently, every
zero-distance class containing a client must contain a candidate, and
there are at most \(k\) such client classes. These conditions are necessary
and sufficient in the original metric; when they hold, open those
locations. Thus the positive-optimum analysis of the reduction does not
require division by zero. Co-located pieces occur only inside our
pseudo-algorithm and do not change this original-metric test.

We also make explicit a rational implementation of the deletion threshold
in the proof of Li--Svensson's reduction~\cite[Section~2, Algorithm~2]{LiSvensson2016}.
That proof analyzes a threshold determined by the unknown optimum.
Its numerical value need not be computed. Suppose a call returns a set
\(T\) with \(k<|T|\le k+c\). For every set \(T'\subseteq T\) obtainable
by at most \(c\) deletions with \(|T'|>k\), and every \(i\in T'\), compute
\[
 \Delta(T',i)=\cost(T'\setminus\{i\})-\cost(T').
\]
These numbers are nonnegative rationals of polynomial bit length.
There are \(n^{O(c)}\) such numbers. Run the reduction's deletion phase
for each distinct number as threshold, and for zero, with a fixed order
for choosing deletions; follow each run by its prescribed enumerative
completion and retain the cheapest feasible result.

To justify this replacement, let \(B^*\ge0\) be the threshold used in
the reduction's analysis. Choose the largest enumerated value at most
\(B^*\), including equality, or zero if there is none. Every test of the form
\(\Delta(T',i)\le B^*\) has the same answer at this representative
threshold. Induction through the at most \(c\) deletions therefore
reproduces the ideal deletion history, including its fixed tie choices.
Its output is present among the enumerated outputs. The approximation
analysis still uses \(B^*\); the algorithm uses only exact rational
comparisons. This adds \(n^{O(c)}\) overhead to the stated reduction.
All candidate enumeration, distance comparisons and cost evaluations
likewise use polynomial-bit rational arithmetic.

Finally, the approximation parameter supplied to the reduction can be a
fixed rational upper bound on our ratio, with its excess charged to the
allowed accuracy. The auxiliary positive constants in its proof may then
be chosen rational as well. These observations preserve both determinism
and polynomial binary-input time; they require no oracle for the optimum.

\begin{proof}[Proof of Theorem~\ref{thm:main}]
Fix \(\varepsilon>0\). Apply Theorem~\ref{thm:pseudo} with
loss \(\varepsilon/4\), obtaining a deterministic
\(c\)-additive approximation, where \(c\) depends only on this fixed
loss. Choose a fixed rational number \(\alpha\) with
\[
 1+2/e+\varepsilon/4<\alpha<1+2/e+\varepsilon/2.
\]
The pseudo-algorithm is a \(c\)-additive \(\alpha\)-approximation.
The additive-to-ordinary reduction of Theorem~\ref{thm:li-svensson},
with a fixed positive rational loss at most \(\varepsilon/2\) and the
implementation in Section~\ref{sec:reduction-implementation},
gives a deterministic polynomial-time
algorithm opening at most \(k\) original facilities and having
cost at most
\[
 \left(1+\frac2e+\varepsilon\right)\OPT_k.
\]
Copies can be merged back to their original facilities as above.
If exactly \(k\) facilities are required, add unused facilities
from \(F\), which is possible because \(k\le|F|\) and can only
decrease the cost.

\end{proof}

Combining Theorem~\ref{thm:main} with the hardness bound recalled in
Section~\ref{sec:introduction} proves Corollary~\ref{cor:threshold}:
under \(P\ne NP\), the infimum of deterministic polynomial-time
approximation factors in the specified candidate-facility model is
\(1+2/e\). The upper bound holds for every fixed positive loss and
does not require attainment of the endpoint.
